\documentclass[a4j, 12pt]{jarticle}
\usepackage[a4paper, left=15mm, right=15mm, top=15mm, bottom=15mm]{geometry}

\usepackage{graphicx}
\usepackage{titling}
\usepackage{amsmath}
\usepackage{float}
\usepackage{listings}
\usepackage{booktabs}
\usepackage{siunitx}
\usepackage{makeidx}
\usepackage{fancybox}
\usepackage{ascmac}

\usepackage{boxedminipage}
\renewcommand{\lstlistingname}{コード}

\graphicspath{{./assets/}}
\title{様々なソートアルゴリズム}
\author{クリスティアン　ハルジュノ\thanks{釧路工業高等専門学校情報工学科5年情報21番}}
\date{\today}
\begin{document}
\begin{titlepage}
  \centering
  \vspace*{3cm}

  {\fontsize{25pt}{36pt}\selectfont\bfseries　パーセプトロン} \\[1cm]  % Main title
  {\Large 釧路工業高等専門学校情報工学科5年} \\[1cm]
  {\Large 人工知能第二課題} \\[1cm]
  {\large 氏名：クリスティアン・ハルジュノ} \\[1cm]
  {\large 出席番号：21} \\[1cm]
  {\large 学籍番号：234071} \\[5cm]

  {\Large 提出日：2025年7月7日} \\

  \vfill

\end{titlepage}
\tableofcontents
\newpage
\setcounter{page}{1}
\section{課題2: 教科書22ページによる\texttt{percep\_main.py}の実行結果}
教科書の22ページによって, 表\ref{page22datasetbefore}ようなデータセットが用意された. 
\begin{table}[H]
  \centering
  \caption{教科書の22ページが用意された学習データ}\label{page22datasetbefore}
  \begin{tabular}{c r}
    \hline
    クラス & 特徴量 $x_1$ \\
    \hline
    $\omega_1$ &  1.2 \\
    $\omega_1$ &  0.2 \\
    $\omega_1$ & -0.2 \\
    $\omega_2$ & -0.5 \\
    $\omega_2$ & -1.0 \\
    $\omega_2$ & -1.5 \\
    \hline
  \end{tabular}
\end{table}

用意されたパーセプトロンのプログラムは2次元特徴空間対象に作られた. しかし, 教科書の２２ページによる学習データには1次元特徴空間であった. その理由で, 表\ref{page22datasetbefore}の学習データを表\ref{page22dataset}のように修正し, ダミー特徴を追加された. また, プログラムが対象できるクラスの数字が制限されているので, $\omega_1$を$0$にして, $\omega_2$を$1$にする. 
\begin{table}[H]
  \centering
  \caption{修正した学習データの一覧(p22.dat)}\label{page22dataset}
  \begin{tabular}{c r r}
    \hline
    クラス & 特徴量$x_1$ & ダミー特徴量 \\
    \hline
    0 &  1.2 &  1.2 \\
    0 &  0.2 &  0.2 \\
    0 & -0.2 & -0.2 \\
    1 & -0.5 & -0.5 \\
    1 & -1.0 & -1.0 \\
    1 & -1.5 & -1.5 \\
    \hline
  \end{tabular}
\end{table}
また, パーセプトロンの訓練を行う時に使用する検証用データも手動で作った.
\begin{table}[H]
  \centering
  \caption{テストデータの一覧(test.dat)}
  \begin{tabular}{c r r}
    \hline
    クラス & 特徴量 $x_1$ & ダミー特徴量 \\
    \hline
    0 &  1.1 &  1.1\\
    0 &  0.5 &  0.5\\
    0 &  0.1 &  0.1\\
    1 & -0.6 & -0.6\\
    1 & -0.8 & -0.8\\
    1 & -1.1 & -1.1\\
    \hline
  \end{tabular}
\end{table}
本実験を与えられたコードそのままで実行すると, パーセプトロンが期待される収束領域の内側ではなく, わずかに外側に位置するという異常な挙動が発生した. この原因は不要な重み $\omega_2$ が含まれていたことである. たとえ $\omega_2$ を0に設定しても, プログラムが重みを含み続けたため, 実験結果が正しくなくなった. この問題を解決するため, 以下のようにコードを修正した. 
\begin{verbatim}
  #変更前
  def to_aug_vector(x):
    dim = x.shape[0]
    aug_x = np.empty([dim + 1, ], dtype=float)
    aug_x[0] = 1.0
    for j in range(0, dim):
        aug_x[j + 1] = x[j]
    return aug_x
  ＃変更後
  def to_aug_vector(x):
    return np.array([1.0, x[0], 0.0], dtype=float)
\end{verbatim}
これらの変更により, $\omega_2$ を常に0に保ち, 他の重みに影響を与えないようにした. 

\subsection{実行方法}
プログラムの実行を容易にするために, 必要なコマンドを1つのコマンドで実行できるよう, Makefile にまとめました. Makefile の内容は次のコードに記載している. 
\begin{verbatim}
  .PHONY: run
  run:
    python3 percep_main.py p22.dat test.dat
      1.2 -7 2 0 weights-test-1.dat logs-test-1.dat
    python3 percep_main.py p22.dat test.dat
      3.6 -7 2 0 weights-test-2.dat logs-test-2.dat
      python3 percep_main.py p22.dat test.dat
    2.0 11 5 0 weights-test-3.dat logs-test-3.dat
    python3 plot-log2d.py logs-test-1.dat 
      logs-test-2.dat logs-test-3.dat output.png
\end{verbatim}
また, 実行する際には次のコマンドのようになる. 
\begin{verbatim}
  make
\end{verbatim}

\subsection{実行結果}
\subsubsection{$\rho=1.2$ $\omega_0=-7$ $\omega_1=2$}
$\rho=1.2$を決定し, 重みベクトル$\omega_0$と$\omega_1$の初期値をそれぞれ$-7$と$2$とする. また, 用意されたプログラムが対象となる重みは三つであるため, $\omega_2$を0にした. \\
プログラムを実行した結果, 表\ref{p12finalweights}に表されている重みになる. 
\begin{table}[H]
  \centering
  \caption{最終的に学習された重みベクトル}\label{p12finalweights}
  \begin{tabular}{l r}
    \hline
    要素 & 値 \\
    \hline
    $w_0$（バイアス項） & $1.400$ \\
    $w_1$（特徴量 $x_1$ の重み） & $3.440$ \\
    $w_2$（ダミー特徴量の重み） & $0.000$ \\
    \hline
  \end{tabular}
\end{table}
表\ref{p12logs}の実行ログから, 本プログラムは7回の反復で学習を完了し, 先に示した表\ref{p12finalweights}に記載された最終的な重みベクトルに到達したことが確認できる. 
\begin{table}[H]
  \centering
  \caption{実行ログにおける重みベクトルの推移（7\,iteration）}\label{p12logs}
  \begin{tabular}{r|rrr}
  \hline
  Step & $w_0$ & $w_1$ & $w_2$ \\
  \hline
  0 & -7.000000 & 2.000000 & 0.000000 \\
  1 & -5.800000 & 3.440000 & 0.000000 \\
  2 & -4.600000 & 3.680000 & 0.000000 \\
  3 & -3.400000 & 3.440000 & 0.000000 \\
  4 & -2.200000 & 3.680000 & 0.000000 \\
  5 & -1.000000 & 3.440000 & 0.000000 \\
  6 &  0.200000 & 3.680000 & 0.000000 \\
  7 &  1.400000 & 3.440000 & 0.000000 \\
  \hline
  \end{tabular}
\end{table}

\subsubsection{$\rho=3.6$ $\omega_0=-7$ $\omega_1=2$}
$\rho=3.6$を決定し, 重みベクトル$\omega_0$と$\omega_1$の初期値をそれぞれ$-7$と$2$とする. また, 用意されたプログラムが対象となる重みは三つであるため, $\omega_2$を0にした. \\
プログラムを実行した結果, 表\ref{p36finalweights}に表されている重みになる. 
\begin{table}[H]
  \centering
  \caption{最終的に学習された重みベクトル}\label{p36finalweights}
  \begin{tabular}{l r}
    \hline
    要素 & 値 \\
    \hline
    $w_0$（バイアス項） & $3.800$ \\
    $w_1$（特徴量 $x_1$ の重み） & $8.480$ \\
    $w_2$（ダミー特徴量の重み）   & $0.000$ \\
    \hline
  \end{tabular}
\end{table}
表\ref{p36logs}の実行ログから, 本プログラムは8回の反復で学習を完了し, 先に示した表\ref{p36finalweights}に記載された最終的な重みベクトルに到達したことが確認できる. しかし, 1.2.1節の前回の実験と比較すると, 学習率（または$\rho$）を$3.6$に大きく設定した場合, パーセプトロンは収束領域を$4$回通り越した後, ようやくその内部に収まった. 
\begin{table}[H]
  \centering
  \caption{実行ログにおける重みベクトルの推移（7\,iteration）}\label{p36logs}
  \begin{tabular}{r|rrr}
    \hline
    Step & $w_0$ & $w_1$ & $w_2$ \\
    \hline
     0 & -7.000000 & 2.000000 & 0.000000 \\
     1 & -3.400000 & 6.320000 & 0.000000 \\
     2 &  0.200000 & 7.040000 & 0.000000 \\
     3 &  3.800000 & 6.320000 & 0.000000 \\
     4 &  0.200000 & 8.120000 & 0.000000 \\
     5 &  3.800000 & 7.400000 & 0.000000 \\
     6 &  0.200000 & 9.200000 & 0.000000 \\
     7 &  3.800000 & 8.480000 & 0.000000 \\
    \hline
  \end{tabular}
\end{table}

\subsubsection{$\rho=2.0$ $\omega_0=11$ $\omega_1=5$}
$\rho=2.0$を決定し, 重みベクトル$\omega_0$と$\omega_1$の初期値をそれぞれ$11$と$5$とする. また, 用意されたプログラムが対象となる重みは三つであるため, $\omega_2$を0にした. \\
プログラムを実行した結果, 表\ref{p20finalweights}に表されている重みになる. 
\begin{table}[H]
  \centering
  \caption{最終的に学習された重みベクトル}\label{p20finalweights}
  \begin{tabular}{l r}
    \hline
    要素 & 値 \\
    \hline
    $w_0$（バイアス項） & $3.000$ \\
    $w_1$（特徴量 $x_1$ の重み） & $10.000$ \\
    $w_2$（ダミー特徴量の重み）   & $0.000$ \\
    \hline
  \end{tabular}
\end{table}
表\ref{p20logs}の実行ログから, 本プログラムは5回の反復だけで学習を完了し, 先に示した表\ref{p20finalweights}に記載された最終的な重みベクトルに到達したことが確認できる. 本実験は, 学習率（または$\rho$）をそれぞれ$1.2$および$3.6$とした1.2.1節および1.2.2節の実験と比較して, 中間的な妥当な結果を示している. 
\begin{table}[H]
  \centering
  \caption{実行ログにおける重みベクトルの推移（4\,iteration）}\label{p20logs}
  \begin{tabular}{r|rrr}
    \hline
    Step & $w_0$ & $w_1$ & $w_2$ \\
    \hline
     0 & 11.000000 & 5.000000 & 0.000000 \\
     1 &  9.000000 & 6.000000 & 0.000000 \\
     2 &  7.000000 & 8.000000 & 0.000000 \\
     3 &  5.000000 & 9.000000 & 0.000000 \\
     4 &  3.000000 & 10.000000 & 0.000000 \\
    \hline
  \end{tabular}
\end{table}
\subsection{考察}
学習率（ρ）を1.2および3.6に設定し, 同じ初期値 $\omega_0 = -7$, $\omega_1 = 2$ から実行した結果を比較すると, パーセプトロンの挙動に違いが見られた. ρ を 1.2 に設定した場合, パーセプトロンは7回の反復で収束領域に到達した. 一方で, ρ を 3.6 に設定した場合も7回の反復で最終的に収束領域に到達したが, 途中で収束領域を複数回通り過ぎており, 収束点にたどり着くまでに安定性を欠いていたことがわかる. これは, 学習率が高すぎることを示しており, より効率的な学習のためには調整が必要であることを示唆している. 
次に, ρ = 1.2（初期値 $\omega_0 = -7$, $\omega_1 = 2$）の結果と, ρ = 2.0（初期値 $\omega_0 = 11$, $\omega_1 = 5$）の結果を比較するため, まずそれぞれの初期点から収束領域の最も近い境界までの距離を計算する. 

式\ref{distanceequation} に示した, 点と直線との距離を求める式を考慮すると, 
\begin{equation}\label{distanceequation}
  d = \frac{|Ax_0 + By_0 + C|}{\sqrt{A^2 + B^2}}
\end{equation}
最も近い収束領域の境界線は, 傾き 0.2 の直線である. $(\omega_0, \omega_1) = (-7, 2)$ からこの直線までの距離を計算する. 
\begin{equation}
  \begin{aligned}
  d &= \frac{\left|Ax_0 + By_0 + C\right|}{\sqrt{A^2 + B^2}} \\
  &= \frac{\left|0.2 \cdot 2 + (-1) \cdot (-7) + 0\right|}{\sqrt{(0.2)^2 + (-1)^2}} \\
  &= \frac{\left|0.4 + 7\right|}{\sqrt{0.04 + 1}} \\
  &= \frac{7.4}{\sqrt{1.04}} \\
  &\approx \frac{7.4}{1.0198} \\
  &\approx 7.26
\end{aligned}
\end{equation}
この計算から, $(\omega_0, \omega_1) = (-7, 2)$ から傾き 0.2 の直線までの距離は $7.26$ であることがわかる. 
次に, $(\omega_0, \omega_1) = (11, 5)$ から傾き 0.5 の直線までの距離を計算する. 
\begin{equation}
  \begin{aligned}
  d &= \frac{|A x_0 + B y_0 + C|}{\sqrt{A^2 + B^2}} \\
  &= \frac{|0.5 \times 5 + (-1) \times 11 + 0|}{\sqrt{(0.5)^2 + (-1)^2}} \\
  &= \frac{|2.5 - 11|}{\sqrt{0.25 + 1}} \\
  &= \frac{8.5}{\sqrt{1.25}} \\
  &= \frac{8.5}{1.1180} \\
  &\approx 7.61
\end{aligned}
\end{equation}
$(\omega_0, \omega_1) = (11, 5)$ の点は, $(\omega_0, \omega_1) = (-7, 2)$ の点よりも最も近い境界線から距離が $7.61 - 7.26 = 0.35$ だけ大きいため, 学習率（ρ）を 2.0 と高く設定しても, パーセプトロンはより短い反復回数で収束領域に到達できる. 
\section{課題3: 課題1によるパーセプトロンの図}
\begin{figure}[H]
  \centering
  \caption{課題1パーセプトロンの図}
  \includegraphics[width=\linewidth]{output.eps}
\end{figure}

\section{課題4: 自作学習データによるパーセプトロン}
自作学習データを作成するため, 複数の条件を確定された. その条件は次のようになる. 

\begin{itembox}[l]{学習パターンの条件}
  \begin{itemize}
    \item $x_{1}$ の範囲：$0 \sim 5$
    \item $x_{2}$ の範囲：$0 \sim 5$
    \item 学習パターン数：各クラス 6 個以上
    \item 線形分離可能な内範囲で，学習アルゴリズムに対してなるべく厳しく
    \item 予想される決定境界が横軸・縦軸と平行にならないように (2特徴を活かすため)
  \end{itemize}
\end{itembox}
その条件に基づいて, 以下のような学習データを作成した.
\begin{table}[H]
  \centering
  \caption{作成した学習データの一覧(mydata.dat)}\label{mydata_table}
  \begin{tabular}{c r r}
    \hline
    クラス & 特徴量$x_1$ & 特徴量$x_2$ \\
    \hline
    $\omega_1$ &  0.5 &  3.0 \\
    $\omega_1$ &  1.0 &  2.5 \\
    $\omega_1$ &  0.8 &  3.1 \\
    $\omega_1$ &  1.2 &  2.6 \\
    $\omega_1$ &  1.5 &  2.0 \\
    $\omega_1$ &  0.6 &  2.8 \\
    $\omega_2$ &  3.0 &  2.0 \\
    $\omega_2$ &  4.0 &  1.0 \\
    $\omega_2$ &  3.5 &  2.0 \\
    $\omega_2$ &  4.5 &  0.5 \\
    $\omega_2$ &  2.5 &  2.2 \\
    $\omega_2$ &  3.2 &  1.5 \\
    \hline
  \end{tabular}
\end{table}
次, この学習データをプロットした結果を図\ref{selfdataplot}に示す. 図から, 2つのクラスが線形分離可能であることが確認できる. 
\begin{figure}[H]
  \centering
  \includegraphics[width=0.5\linewidth]{plotted_self_data.eps}
  \caption{自作学習データのプロット}\label{selfdataplot}
\end{figure}

また, パーセプトロンを訓練するための検証用データも手動で作成した. 検証用データは, 学習データと同じクラス分布を持つが, 異なる特徴量を持つ. 
\begin{table}[H]
  \centering
  \caption{テストデータの一覧(mytest.dat)}\label{mytest_table}
  \begin{tabular}{c r r}
    \hline
    クラス & 特徴量$x_1$ & 特徴量$x_2$ \\
    \hline
    1 &  0.7 &  3.4 \\
    1 &  0.9 &  3.0 \\
    1 &  1.1 &  2.8 \\
    1 &  1.3 &  2.5 \\
    1 &  1.6 &  1.9 \\
    1 &  0.4 &  3.2 \\
    0 &  3.4 &  1.1 \\
    0 &  3.8 &  1.2 \\
    0 &  4.2 &  0.6 \\
    0 &  3.6 &  1.4 \\
    0 &  2.8 &  2.0 \\
    0 &  3.1 &  1.3 \\
    \hline
  \end{tabular}
\end{table}

\subsection{実行方法}
自作学習データを使用してパーセプトロンを実行するために, 以下のようなMakefileを作成した. このMakefileは, 自作学習データを使用してパーセプトロンを訓練し, 結果をプロットするためのコマンドをまとめている.この実験では, 3つの異なる学習率（$\rho$）を使用してパーセプトロンを訓練し, それぞれの結果を比較する. 以下はMakefileの内容である. 
\begin{verbatim}
	python3 percep_main.py mydata.dat mytest.dat
    1.0 7.5 -6.2 3.3 weights-self-1.dat logs-self-1.dat
	python3 percep_main.py mydata.dat mytest.dat
    2.0 -9.1 4.4 -7.8 weights-self-2.dat logs-self-2.dat
	python3 percep_main.py mydata.dat mytest.dat
    10.0 11.0 2.1 -3.5 weights-self-3.dat logs-self-3.dat
	python3 plot-log3d.py logs-self-1.dat logs-self-2.dat
    logs-self-3.dat output-self.png
\end{verbatim}
\subsection{実行結果}
\subsubsection{$\rho=1.0$ $\omega_0=7.5$ $\omega_1=-6.2$ $\omega_2=3.3$}
$\rho=1.0$を決定し, 重みベクトル$\omega_0$と$\omega_1$の初期値をそれぞれ$7.5$と$-6.2$とする. また, 用意されたプログラムが対象となる重みは三つであるため, $\omega_2$を$3.3$にした. \\
プログラムを実行した結果, 表\ref{self1finalweights}に表されている重みになる. 
\begin{table}[H]
  \centering
  \caption{最終的に学習された重みベクトル}\label{self1finalweights}
  \begin{tabular}{l r}
    \hline
    要素 & 値 \\
    \hline
    $w_0$（バイアス項） & $5.500$ \\
    $w_1$（特徴量 $x_1$ の重み） & $3.800$ \\
    $w_2$（特徴量 $x_2$ の重み） & $-6.200$ \\
    \hline
  \end{tabular}
\end{table}
表\ref{self1logs}の実行ログから, 本プログラムは12回の反復で学習を完了し, 先に示した表\ref{self1finalweights}に記載された最終的な重みベクトルに到達したことが確認できる.
\begin{table}[H]
  \centering
  \caption{実行ログにおける重みベクトルの推移（13\,iteration）}\label{self1logs}
  \begin{tabular}{r|rrr}
  \hline
  Step & $w_0$ & $w_1$ & $w_2$ \\
  \hline
   0 &  7.500000 & -6.200000 &  3.300000 \\
   1 &  6.500000 & -6.700000 &  0.300000 \\
   2 &  5.500000 & -7.700000 & -2.200000 \\
   3 &  6.500000 & -4.700000 & -0.200000 \\
   4 &  7.500000 & -0.700000 &  0.800000 \\
   5 &  6.500000 & -1.200000 & -2.200000 \\
   6 &  5.500000 & -2.700000 & -4.200000 \\
   7 &  6.500000 &  0.300000 & -2.200000 \\
   8 &  5.500000 & -0.200000 & -5.200000 \\
   9 &  6.500000 &  2.800000 & -3.200000 \\
  10 &  5.500000 &  1.800000 & -5.700000 \\
  11 &  6.500000 &  4.800000 & -3.700000 \\
  12 &  5.500000 &  3.800000 & -6.200000 \\
  \hline
  \end{tabular}
\end{table}
\subsubsection{$\rho=2.0$ $\omega_0=-9.1$ $\omega_1=4.4$ $\omega_2=-7.8$}
$\rho=2.0$を決定し, 重みベクトル$\omega_0$と$\omega_1$の初期値をそれぞれ$-9.1$と$4.4$とする. また, 用意されたプログラムが対象となる重みは三つであるため, $\omega_2$を$-7.8$にした. \\
プログラムを実行した結果, 表\ref{self2finalweights}に表されている重みになる. 
\begin{table}[H]
  \centering
  \caption{最終的に学習された重みベクトル}\label{self2finalweights}
  \begin{tabular}{l r}
    \hline
    要素 & 値 \\
    \hline
    $w_0$（バイアス項） & $-9.100$ \\
    $w_1$（特徴量 $x_1$ の重み） & $13.000$ \\
    $w_2$（特徴量 $x_2$ の重み）） & $-8.600$ \\
    \hline
  \end{tabular}
\end{table}
表\ref{self2logs}の実行ログから, 本プログラムは6回の反復で学習を完了し, 先に示した表\ref{self2finalweights}に記載された最終的な重みベクトルに到達したことが確認できる.
\begin{table}[H]
  \centering
  \caption{実行ログにおける重みベクトルの推移（7\,iteration）}\label{self2logs}
  \begin{tabular}{r|rrr}
  \hline
  Step & $w_0$ & $w_1$ & $w_2$ \\
  \hline
   0 & -9.100000 &  4.400000 & -7.800000 \\
   1 & -7.100000 & 10.400000 & -3.800000 \\
   2 & -9.100000 &  7.400000 & -7.800000 \\
   3 & -7.100000 & 13.400000 & -3.800000 \\
   4 & -9.100000 & 10.400000 & -7.800000 \\
   5 & -7.100000 & 15.400000 & -3.400000 \\
   6 & -9.100000 & 13.000000 & -8.600000 \\
  \hline
  \end{tabular}
\end{table}
\subsubsection{$\rho=10.0$ $\omega_0=11.0$ $\omega_1=2.1$ $\omega_2=-3.5$}
$\rho=10.0$を決定し, 重みベクトル$\omega_0$と$\omega_1$の初期値をそれぞれ$11.0$と$2.1$とする. また, 用意されたプログラムが対象となる重みは三つであるため, $\omega_2$を$-3.5$にした. \\
プログラムを実行した結果, 表\ref{self3finalweights}に表されている重みになる. 
\begin{table}[H]
  \centering
  \caption{最終的に学習された重みベクトル}\label{self3finalweights}
  \begin{tabular}{l r}
    \hline
    要素 & 値 \\
    \hline
    $w_0$（バイアス項） & $1.000$ \\
    $w_1$（特徴量 $x_1$ の重み） & $52.100$ \\
    $w_2$（特徴量 $x_2$ の重み）& $-46.500$ \\
    \hline
  \end{tabular}
\end{table}
表\ref{self3logs}の実行ログから, 本プログラムは4回の反復で学習を完了し, 先に示した表\ref{self3finalweights}に記載された最終的な重みベクトルに到達したことが確認できる.
\begin{table}[H]
  \centering
  \caption{実行ログにおける重みベクトルの推移（8\,iteration）}\label{self3logs}
  \begin{tabular}{r|rrr}
  \hline
  Step & $w_0$ & $w_1$ & $w_2$ \\
  \hline
   0 & 11.100000 &  2.100000 &  -3.500000 \\
   1 &  1.000000 & -2.900000 & -33.500000 \\
   2 & 11.100000 & 27.100000 & -13.500000 \\
   3 &  1.000000 & 17.100000 & -38.500000 \\
   4 & 11.100000 & 47.100000 & -18.500000 \\
   5 &  1.000000 & 37.100000 & -43.500000 \\
   6 & 11.100000 & 62.100000 & -21.500000 \\
   7 &  1.000000 & 52.100000 & -46.500000 \\
  \hline
  \end{tabular}
\end{table}
\subsection{考察}
収束領域が非常に狭いため、学習率を低く設定しても収束に至らない状況が繰り返されることは、さらに小さな学習率、例えば $0.01$ や $0.05$ の設定が必要である可能性を示している。学習率を小さくすることにより、収束領域内でオーバーシュートせずに停止するなど、学習の安定性は大幅に向上する可能性がある。しかしその一方で、学習率が小さすぎると、最終的な収束点に到達するまでに必要な反復回数は大幅に増加する。
実際に、学習率 $\rho = 0.05$、初期重みを $\omega_0 = 10.0, \omega_1 = 11.5, \omega_2 = 2.0$ に設定してシミュレーションを実行したところ、表 \ref{smalllearningratelogs} に示すように、収束点に到達するまでの反復回数は顕著に増加した。
\begin{table}[H]
  \centering
  \caption{実行ログにおける重みベクトルの推移（抜粋）}\label{smalllearningratelogs}
  \begin{tabular}{r|rrr}
  \hline
  Step & $w_0$ & $w_1$ & $w_2$ \\
  \hline
   0 & 10.000000 & 11.500000 &  2.000000 \\
   1 &  9.950000 & 11.474999 &  1.850000 \\
   2 &  9.900000 & 11.424999 &  1.725000 \\
   3 &  9.850000 & 11.384999 &  1.570000 \\
   \vdots & \vdots & \vdots & \vdots \\
  83 &  5.850000 &  6.885000 & -8.260000 \\
  \hline
  \end{tabular}
\end{table}

最適な学習率を一意に決定するアルゴリズムや手法は存在しない。そのため、現実的なアプローチとしては、繰り返しの試行錯誤によって最適値を探索することが必要となる。また、学習率を徐々に減衰させる「学習率減衰（learning rate decay）」を導入すれば、序盤は大きなステップで収束領域に接近し、終盤は細かく調整することでオーバーシュートを防ぐことが可能となる。
\section{課題5:　グラフに推測した解領域}
\begin{figure}[H]
  \centering
  \includegraphics[width=\linewidth]{cropped1.eps}
  \caption{推測した解領域のビュー1}
\end{figure}
\begin{figure}[H]
  \centering
  \includegraphics[width=0.6\linewidth]{cropped2.eps}
  \caption{推測した解領域のビュー2}
\end{figure}
\begin{figure}[H]
  \centering
  \includegraphics[width=0.6\linewidth]{cropped3.eps}
  \caption{推測した解領域のビュー3}
\end{figure}
\begin{figure}[H]
  \centering
  \includegraphics[width=0.6\linewidth]{cropped4.eps}
  \caption{推測した解領域のビュー4}
\end{figure}
\section{課題6:　解領域と学習パターンをプロット}
前回行った実験には二つの特徴（2次元特徴量）を持つ学習パターンが与えられた。これらの特徴量は、2次元空間上で線形分離可能な領域を形成している。境界線の位置を計算するために、式\ref{decisionboundary}を使用する。ここで、$w_0$はバイアス項、$w_1$は特徴量$x_1$の重み、$w_2$は特徴量$x_2$の重みを表す。
\begin{equation}\label{decisionboundary}
  w_0 + w_1 x_1 + w_2 x_2 = 0
  \quad \Rightarrow \quad
  x_2 = -\frac{w_1}{w_2} x_1 - \frac{w_0}{w_2}
\end{equation}
表\ref{self1finalweights}、\ref{self2finalweights}、および\ref{self3finalweights}に示した最終重みを基に、式\ref{decisionboundary}を適用する。結果は表\ref{boundarycalculation} に示す通りである。
\begin{table}[H]
  \centering
  \caption{最終的な重みベクトルから計算した傾きと切片}\label{boundarycalculation}
  \begin{tabular}{c r r r r r}
    \hline
    ケース & $w_0$ & $w_1$ & $w_2$ & 傾き $-\frac{w_1}{w_2}$ & 切片 $-\frac{w_0}{w_2}$ \\
    \hline
    1 & 5.5 & 3.8 & -6.2 & 0.613 & 0.887 \\
    2 & -9.1 & 13.0 & -8.6 & 1.512 & -1.058 \\
    3 & 1.0 & 52.1 & -46.5 & 1.120 & 0.022 \\
    \hline
  \end{tabular}
\end{table}
ケース1は3.2.1章によって、最終的な重みベクトルが $w_0 = 5.5$, $w_1 = 3.8$, $w_2 = -6.2$ の場合である。この場合、傾きは約 $0.613$、切片は約 $0.887$ となる。ケース2では3.2.2章によって、$w_0 = -9.1$, $w_1 = 13.0$, $w_2 = -8.6$ の場合で、傾きは約 $1.512$、切片は約 $-1.058$ となる。ケース3では3.2.3章によって、では、$w_0 = 1.0$, $w_1 = 52.1$, $w_2 = -46.5$ の場合で、傾きは約 $1.120$、切片は約 $0.022$ となる。全ケースをプロットすると\ref{allcasesplot}のようになる。ここでは、各ケースの解領域と学習パターンをプロットしている。
\begin{figure}[H]
  \centering
  \includegraphics[width=0.72\linewidth]{all-case-plot.eps}
  \caption{三つのケースの解領域と学習パターンのプロット}\label{allcasesplot}
\end{figure}
\subsection{ケース1の解領域と学習パターンのプロット}
\begin{table}[H]
  \centering
  \caption{ケース1の最終的な重みベクトルから計算した傾きと切片}
  \begin{tabular}{c r r r r r}
    \hline
    ケース & $w_0$ & $w_1$ & $w_2$ & 傾き $-\frac{w_1}{w_2}$ & 切片 $-\frac{w_0}{w_2}$ \\
    \hline
    1 & 5.5 & 3.8 & -6.2 & 0.613 & 0.887 \\
    \hline
  \end{tabular}
\end{table}
\begin{figure}[H]
  \centering
  \includegraphics[width=0.7\linewidth]{case1-plot.eps}
  \caption{ケース1の解領域と学習パターンのプロット}\label{case1plot}
\end{figure}
\newpage
\subsection{ケース2の解領域と学習パターンのプロット}
\begin{table}[H]
  \centering
  \caption{ケース2の最終的な重みベクトルから計算した傾きと切片}
  \begin{tabular}{c r r r r r}
    \hline
    ケース & $w_0$ & $w_1$ & $w_2$ & 傾き $-\frac{w_1}{w_2}$ & 切片 $-\frac{w_0}{w_2}$ \\
    \hline
    2 & -9.1 & 13.0 & -8.6 & 1.512 & -1.058 \\
    \hline
  \end{tabular}
\end{table}
\begin{figure}[H]
  \centering
  \includegraphics[width=0.7\linewidth]{case2-plot.eps}
  \caption{ケース2の解領域と学習パターンのプロット
  }\label{case2plot}
\end{figure}
\newpage
\subsection{ケース3の解領域と学習パターンのプロット}
\begin{table}[H]
  \centering
  \caption{ケース3の最終的な重みベクトルから計算した傾きと切片}
  \begin{tabular}{c r r r r r}
    \hline
    ケース & $w_0$ & $w_1$ & $w_2$ & 傾き $-\frac{w_1}{w_2}$ & 切片 $-\frac{w_0}{w_2}$ \\
    \hline
    3 & 1.0 & 52.1 & -46.5 & 1.120 & 0.022 \\
    \hline
  \end{tabular}
\end{table}
\begin{figure}[H]
  \centering
  \includegraphics[width=0.7\linewidth]{case3-plot.eps}
  \caption{ケース3の解領域と学習パターンのプロット
  }\label{case3plot}
\end{figure}
\section{まとめ}
この実験では、パーセプトロンの学習における学習率の調整が与える影響を検証した。学習率を下げるほど学習は安定するが、反復回数が著しく増加する。一方で、学習率を上げすぎると収束領域を越えてしまい、パーセプトロンが前後に揺れ動く状態となる。高い学習率は反復回数を減らし学習を高速化するが、学習の安定性は著しく低下する。また、学習データの分布や初期重みの設定も学習の収束に大きく影響することが分かった。適切な学習率と初期値の選択が、効率的かつ安定した学習には不可欠である。さらに、学習率減衰などの工夫を導入することで、初期段階の高速な収束と最終段階の安定性を両立できる可能性が示唆された。
\end{document}
